\chapter{Future Work}
\label{sec:future_work}

\section{Broader Evaluation}
One possible extension is to evaluate additional query collections with clearly
different characteristics. Such an extension would test how sensitive the
comparison is to the initial study setting. It should preserve enough of the
original protocol to make the new evidence interpretable while explaining
any changes required by the additional data.

A broader collection is not automatically a better experiment. Its relevance
judgements, sampling procedure and documentation still determine what can be
learned from it. Before expanding the study, identify the uncertainty that the
new collection is intended to reduce. This makes the extension a targeted
investigation rather than simply a larger repetition of the same procedure.

\section{User-Centred Analysis}
Another extension would examine how people interact with the retrieved results.
An offline score represents selected aspects of retrieval quality, whereas
a user study can investigate task completion, interpretation or effort under
a specified interface. The study design would need to define the tasks and
outcomes carefully so that the observations answer a focused question.

The interface, instructions and participant population would become part of
the experimental setting. These elements should be documented alongside the
retrieval configuration. A user study should not be presented as a direct
replacement for an offline comparison; the two approaches can provide evidence
about different aspects of the same system.

% Demonstration-only break.
\newpage
\section{Efficiency and Deployment}
A further direction is to examine the computational cost of the reranking
stage. Relevant measurements might include response time, memory use and the
number of candidates processed. The evaluation should state how these values
were measured and which parts of the pipeline were included. A timing result
without the corresponding execution conditions is difficult to interpret.

The eventual deployment setting may impose a different balance between ranking
quality and computational cost. A small improvement may be useful in one
application but insufficient to justify added delay in another. Establishing
that balance requires a concrete account of the task and its constraints,
rather than a universal claim about the best configuration.

\section{Prioritising the Next Step}
A practical research plan should identify which extension would be most
informative and what resources it requires. The next step could be a targeted
error analysis, an additional collection or a measurement of runtime cost.
Its value depends on which uncertainty currently limits the interpretation
of the study, not simply on how many new components it introduces.

This completes the numbered chapters of the example. Acknowledgements and
references follow as unnumbered chapters, each with a visible page number on
its opening page. The appendix then demonstrates lettered chapter numbering.
These transitions are included so that the same page-style rules can be
checked across the complete thesis structure.
